% 本文件是内容脚手架，不是固定目录。只保留与本题证据链有关的部分，删除不适用项。
% 禁止把示例行、待替换文字、固定问题数或固定“识别—预测—优化”顺序带入终稿。

\section{问题重述与任务拆解}

\subsection{题面信息与口径裁决}

% 用自己的话压缩题目背景，并逐项确认数据、已知条件、硬约束、目标输出与歧义裁决。
% 不复制题面原句，不在本节提前堆模型名、公式或未经验证的结论。
用 2--4 个自然段说明研究对象、任务边界和统一口径。题面存在歧义时，写明采用的解释、依据及其对后续结果的影响。

\subsection{任务映射与成功条件}

% 每个编号问题占一行；按真实问题数增删。方法尚未选定时只写数学任务，不先填算法名。
% “真实依赖”只记录确实传递的数据、参数、结果或约束；相互独立时明确写“无”。
各子问题的输入、数学任务、输出与验收口径见表~\ref{tab:task-map}。该表先定义“怎样算答对”，再进入模型选择，避免求解后临时更换评价标准。

\begin{table}[htbp]
  \centering
  \caption{各子问题的任务映射与成功条件}
  \label{tab:task-map}
  \begin{tabularx}{\textwidth}{>{\centering\arraybackslash}p{0.10\textwidth}XXXX}
    \toprule
    子问题 & 输入与硬约束 & 数学任务 & 输出与成功条件 & 真实依赖 \\
    \midrule
    问题~$i$ & \textbf{待替换：}题面数据、已知条件与约束 & \textbf{待替换：}需要计算、解释或决策的对象 & \textbf{待替换：}交付物、指标、单位与阈值 & \textbf{待替换：}上游/下游接口或“无” \\
    \bottomrule
  \end{tabularx}
\end{table}

表后只解释跨问共享口径、真正的数据流和误差传播；没有依赖的子问题保持独立，不为叙事连贯强造接口。

% 仅当多问存在明显数据流、模型递进、共享验证边界或方法选择关系时保留下一章。
% 简单题直接进入模型假设、主要符号说明和各问题章节，避免与每问导读重复。
\section{总体建模框架与证据链}

\subsection{统一建模思路}

% 用本题对象写清“事实/数据 → 变量与口径 → 基线 → 主模型 → 结果 → 独立检验 → 结论”。
% 复杂多问论文若确需总流程图，须按 technical-roadmap.md 由真实证据生成 F01；不要复制通用空框图。
围绕题目的真实对象概括统一主线，说明各问题怎样共享数据、参数或验证边界，以及每个关键结论如何回链到数据、公式、算法输出和独立检验。

\subsection{复杂度控制与验证安排}

% 先写最简可解释基线；只有同口径验证证明新增模块有收益时，才保留复杂模型。
% 验证方式服从 IID、分组、时间、空间或优化可行性结构，不能机械套用五折交叉验证。
列出基线、主模型和候选改进的保留条件，并在建模前冻结主要评价指标、数据隔离方式、稳健性检查和失败判据。新增复杂度必须对应可量化收益，同时报告计算代价、解释代价和适用边界。

\subsection{共享数据审计与变量口径}

% 仅当多个问题复用同一批数据或变量定义时保留；否则把审计放回对应问题章节。
% 必查：样本量、字段、主键、缺失、异常、重复、时间/空间范围、单位、来源和泄漏风险。
报告全局共享的数据质量结论、变量方向、单位/量纲、变换依据与训练—验证—测试边界。每项处理说明原因和潜在偏差，不只写“完成清洗”或“进行了归一化”。
